\documentclass[11pt,a4paper]{amsart}
\usepackage[T1]{fontenc}
\usepackage{amssymb}
\usepackage[margin=1in]{geometry}
\usepackage[hidelinks]{hyperref}
\hypersetup{pdftitle={A precise independent e-merger outside martingale mixtures}}
\DeclareFontShape{OT1}{cmr}{bx}{n}{<-> cmbx10}{}
\DeclareMathSizes{8}{8}{8}{8}
\DeclareMathSizes{9}{9}{8}{8}
\DeclareMathSizes{10.95}{10.95}{8}{8}
\newtheorem{theorem}{Theorem}
\title{A precise independent e-merger outside martingale mixtures}
\date{}
\begin{document}
\begin{abstract}
A continuous nonmonotone merger of three independent e-values is precise,
yet lies outside the generalized martingale mixtures of Vovk and Wang.
The construction applies the known Ara\'ujo--Feix/Baumeler--Wolf Boolean
selector to the statistical setting. Four rational inputs give a linear
constraint on every adaptive-order wealth mixture and separate the
example by at least \(3/52\) in maximum error on these inputs.
Precision also excludes pointwise domination by such a mixture.
This answers Question 4 of Vovk and Wang.
\end{abstract}
\maketitle

\section{Definitions and result}

An e-variable is a nonnegative random variable \(E\) with
\(\mathbb E E\leq1\). An independent e-merger is a Borel function
\(F:[0,\infty)^K\to[0,\infty)\) such that
\(\mathbb E F(E_1,\ldots,E_K)\leq1\) whenever the \(E_i\) are independent
e-variables. It is precise if the output has mean one whenever all
input means are one. We allow each input to have its own law.

We use the generalized martingale merging class of
Vovk and Wang \cite[\S6, eqs.~(18)--(19)]{vw}.
A strategy starts with capital \(c\in[0,1]\).
At each step it picks an unread index and a fraction \(t\in[0,1]\)
using the revealed history. Reading \(x_i\) multiplies its wealth
by \(1-t+tx_i\). Each input is read exactly once.
A generalized martingale merger \(M\) averages terminal wealth under a
probability law fixed across input vectors. The law may couple the
reading and betting rules. Their Theorem 3 gives validity of every
such mixture. Their Section 7 Question 4 asks whether every precise
independent e-merger has this representation.

\begin{theorem}
There is a continuous precise independent e-merger \(F\) on three
inputs such that no generalized martingale merger dominates it
pointwise. For every such mixture \(M\), the maximum absolute
error from \(F\) on the four inputs
\[
 A=(1/4,1/4,1/4),\quad B_1=(0,1/4,3/4),\quad
 B_2=(3/4,0,1/4),\quad B_3=(1/4,3/4,0)
\]
is at least \(3/52\). The same lower bound holds for pointwise limits
of generalized martingale mergers.
\end{theorem}

\section{Construction and precision}

Let \(r:[0,\infty)\to[0,1]\) be a Borel cutoff and write
\(a_i=r(x_i)\). Define
\[
 w_1=(1-a_2)a_3,\qquad
 w_2=(1-a_3)a_1,\qquad
 w_3=(1-a_1)a_2,\qquad
 F(x)=1+\sum_{i=1}^3 (x_i-1)w_i .
\]
On Boolean inputs the map \(a\mapsto (w_1,w_2,w_3)\) is the Lugano
(Ara\'ujo--Feix/Baumeler--Wolf) process \cite[Example~4 and Table~I]{bw}.
Its explicit Boolean equations also appear in \cite[eq.~(2)]{kb}.
Each weight ignores its own input. Each pair of Boolean flags
requires opposite values of a shared bit, so at most one is active.
More generally, expanding the two products gives
\[
 1-\sum_{i=1}^3 w_i = a_1a_2a_3+(1-a_1)(1-a_2)(1-a_3)\geq0 .
\]
Thus \(F\) is a pointwise convex combination of \(1,x_1,x_2,x_3\).
It is Borel, nonnegative and bounded by \(1+\sum_i x_i\).
If the \(E_i\) are independent, \(E_i\) is independent of
\(w_i=w_i(E_1,E_2,E_3)\). Since \(0\leq w_i\leq1\), the summands
are integrable and
\[
 \mathbb E F(E_1,E_2,E_3)=
  1+\sum_{i=1}^3 (\mathbb E E_i-1)\mathbb E w_i.
\]
This expression is at most one for e-variables and equals one if
all three means are one.

For the continuous example choose
\(r(x)=\min\{1,\max\{0,2x-1/2\}\}\).
It is zero up to \(1/4\), one from \(3/4\) onward, and linear in
between. The hard cutoff \(r(x)=\mathbf{1}_{\{x>1/2\}}\) gives
a second example with the same four witness values. At \(A\) all
weights vanish, so \(F(A)=1\). At \(B_i\) exactly \(w_i=1\)
and \(x_i=0\), so \(F(B_i)=0\).
For either cutoff \(F(0,1/4,1/4)=1>F(0,1/4,3/4)=0\).
In particular these examples lie outside the coordinatewise
increasing restriction of Question 3 in \cite{vw}.

\section{Four inputs rule out every mixture}

For each deterministic strategy, its first index \(i\) and first
fraction \(t\) are independent of the input vector, since there is
no revealed history. Write \(W(x)\) for its terminal wealth.
At \(A\) all later multipliers are at most one. Consequently
\[
 W(A)\leq c(1-3t/4)=c/4+(3/4)c(1-t)\leq1 .
\]
If a mixture had \(M(A)=F(A)=1\), this bound and nonnegativity would
force \(c=1,t=0\) for almost all components. At \(B_i\), a component
reading \(i\) first with no bet still has capital one after seeing
the zero. Its other two inputs are strictly positive, so both
multipliers and terminal wealth are strictly positive. The mixture
requirement \(M(B_i)=0\) rules out positive probability for this
first index. This holds for every \(i\), a contradiction.

A quantitative form avoids the exact-equality assumption.
At \(B_i\), for the same component starting with index \(i\),
the first multiplier equals \(1-t\). The later multipliers
are at least \(1/4\) and \(3/4\), in whichever order they occur,
because a convex combination of one and \(x\leq1\) is at least \(x\).
Hence
\[
 W(B_i)\geq (3/16)c(1-t).
\]
All other terminal values are nonnegative. Together,
\[
 W(A)\leq\tfrac14+4\sum_{j=1}^3W(B_j).
\]
Average this inequality over the fixed strategy law. This keeps the
bound and replaces \(W\) with \(M\). All the evaluations are bounded
by one on these inputs, so they are integrable.
If the maximum error on the four points is \(\delta\), then
\(1-\delta\leq M(A)\leq1/4+12\delta\), giving \(\delta\geq3/52\).
Passing to a pointwise limit preserves the inequality because only
four evaluations appear.

\section{Admissibility and earlier work}

Every precise independent e-merger is admissible among independent
e-mergers by Ramdas and Wang \cite[Prop.~8.6]{rw}.
For completeness, suppose a valid \(G\geq F\) has a strict inequality
at a finite vector \(x\). For each coordinate choose
\(q_i\in(0,1)\) small enough that \(q_i x_i<1\).
Put probability \(q_i\) at \(x_i\) and the rest at
\((1-q_i x_i)/(1-q_i)\). The independent product law has mean-one
marginals and positive mass at \(x\). Under it, \(G\) has expectation
strictly greater than one, contradicting validity. Thus
domination by a generalized martingale merger would imply
representation, already excluded by the four evaluations.

Vovk and Wang's Example 4, credited there to Zhenyuan Zhang,
already gives a two-input nondomination example. Precision is
absent there. For its threshold \(h>1\), a constant-one first input
and a second input taking \(h\) with probability \(1/h\) and zero
otherwise have output mean \(1-1/h\).
The known Lugano causal-game results distinguish the full
multibit process from causally ordered channels.
Here the statistical conclusion uses a scalar merger and the
specific adaptive wealth law of \cite{vw}.
The Boolean rule is prior work. The new question addressed
here is the representation of its precise statistical readout.

\section{Verification and disclosure}

Internally reviewed, not peer reviewed. The rational replay checks
the four evaluations, Boolean and fractional flag sums, product-law
examples, and finite betting arithmetic. The partial Lean 4.30.0
certificate checks an algebraic identity, witness evaluations,
rational pathwise betting facts and the numerical inequality for
the \(3/52\) bound. The expectation, admissibility and representation
arguments are given in the text.

Under human direction, ChatGPT Astra 6.2 (OpenAI) generated the
discovery proof and its research packet. OpenAI Codex performed
the independent mathematical review and developed the partial
Lean certificate. GPT-6.2 Sol (OpenAI) checked the literature,
derived the quantitative bound, and prepared the manuscript,
replay and illustrations. DannyExperiments maintains the repository
and is the preferred citation identity. No AI system is proposed
as an author.

\begin{thebibliography}{9}
\bibitem{vw}
V. Vovk and R. Wang, \emph{Merging sequential e-values via martingales},
Electronic Journal of Statistics \textbf{18} (2024), 1185--1205.
\href{https://doi.org/10.1214/24-EJS2228}{doi:10.1214/24-EJS2228}.
\bibitem{rw}
A. Ramdas and R. Wang, \emph{Hypothesis testing with e-values},
Foundations and Trends in Statistics \textbf{1} (2025), 1--390.
\href{https://doi.org/10.1561/3600000002}{doi:10.1561/3600000002}.
\bibitem{bw}
\"A. Baumeler and S. Wolf,
\emph{The space of logically consistent classical processes without
causal order}, New Journal of Physics \textbf{18} (2016), 013036.
\href{https://doi.org/10.1088/1367-2630/18/1/013036}{doi:10.1088/1367-2630/18/1/013036}.
\bibitem{kb}
R. Kunjwal and \"A. Baumeler,
\emph{Quantum nonlocality without entanglement through multipartite
indefinite causal order}, Physical Review Letters \textbf{131}
(2023), 120201.
\href{https://doi.org/10.1103/PhysRevLett.131.120201}{doi:10.1103/PhysRevLett.131.120201}.
\end{thebibliography}
\end{document}
